\documentclass[11pt]{article}
\usepackage[margin=0.82in]{geometry}
\usepackage{amsmath,amssymb,amsthm}
\usepackage{graphicx}
\usepackage{microtype}
\usepackage{xcolor}
\usepackage{hyperref}
\hypersetup{colorlinks=true,linkcolor=blue!55!black,urlcolor=blue!55!black,citecolor=blue!55!black}
\newtheorem{theorem}{Theorem}
\newtheorem{lemma}{Lemma}
\newcommand{\E}{\mathbb E}
\newcommand{\R}{\mathbb R}
\newcommand{\Var}{\operatorname{Var}}
\title{Sharp negative spherical moments when one coefficient is large\\
a partial answer to arXiv:2206.01026}
\author{Candidate partial-result packet}
\date{13 August 2026}

\begin{document}
\maketitle

\begin{center}
\fbox{\parbox{0.92\linewidth}{\textbf{Status:} candidate partial result likely valid.  The conjectured two-summand constant is proved throughout the source's last negative-moment strip whenever one coefficient carries at least half of the total squared mass.  The complementary all-small regime remains open.}}
\end{center}

\section{Source questions}

Let $\xi_1,\xi_2,\ldots$ be independent random vectors, each uniform on
$S^{d-1}\subset\R^d$.  For $q>-(d-1)$, the source defines $c_d(q)$ as the
best constant such that
\begin{equation}\label{eq:khin}
 \left\|\sum_k a_k\xi_k\right\|_q
 \geq c_d(q)\left(\sum_k a_k^2\right)^{1/2}.
\end{equation}
The two-equal-coefficient candidate is
\[
 c_{d,2}(q)=\left\|\frac{\xi_1+\xi_2}{\sqrt2}\right\|_q.
\]
Chasapis, Singh, and Tkocz state that the following ranges appear open and
conjecture the two-summand value~\cite{CST}:
\[
 c_2(q)=c_{2,2}(q)\quad(-1<q<0),\qquad
 c_3(q)=c_{3,2}(q)\quad(-2<q<-1).
\]

\begin{center}
\includegraphics[width=0.91\textwidth]{figures/open_problem_crop.png}

\smallskip
\emph{The two open ranges in Section 1.1, page 3 of the source PDF.}
\end{center}

\section{Proof intuition}

Write $q=-p$.  Because the exponent is negative, the desired lower bound for
the $L_q$ norm is an \emph{upper} bound for the singular moment
$\E|\sum a_k\xi_k|^{-p}$.  Project onto a largest coefficient $a$ and call
the remaining random sum $Y$.  Conditional on $Y$, the moment is a spherical
Riesz potential evaluated at $|Y|/a$.

In the last strip $d-2<p<d-1$, the source's hypergeometric expansion has
strictly positive nonconstant coefficients.  Consequently this potential
increases until the two vectors have equal lengths and decreases thereafter.
It is therefore bounded by the equal-length value, independently of the
random tail.  If $a$ contains at least half of the squared mass, the remaining
normalization is exactly the conjectured two-summand constant.

\section{The partial theorem}

For $0<p<d-1$, set
\[
 M_{d,2}(p)=\E\left|\frac{\xi_1+\xi_2}{\sqrt2}\right|^{-p}
 =c_{d,2}(-p)^{-p}.
\]

\begin{theorem}[large-coefficient sharp inequality]\label{thm:main}
Let $d\geq2$ and $d-2<p<d-1$.  For every finite real coefficient vector
$a=(a_1,\ldots,a_n)\ne0$ satisfying
\begin{equation}\label{eq:large}
 \max_k a_k^2\geq\frac12\sum_k a_k^2,
\end{equation}
one has
\begin{equation}\label{eq:moment}
 \E\left|\sum_{k=1}^n a_k\xi_k\right|^{-p}
 \leq M_{d,2}(p)\left(\sum_{k=1}^n a_k^2\right)^{-p/2}.
\end{equation}
Equivalently, \eqref{eq:khin} holds with the conjectured constant
$c_{d,2}(-p)$ for every such vector.  Equality in \eqref{eq:moment} occurs
exactly when, up to permutation, signs, and zero coefficients,
$|a_1|=|a_2|>0$.
\end{theorem}

Thus the theorem covers the whole large-coefficient region for both open
low-dimensional ranges, and more generally for the last strip
$-(d-1)<q<-(d-2)$ in every $d\geq2$.

\section{Proof}

\begin{lemma}[spherical potential peaks at matching radii]\label{lem:peak}
For $d-2<p<d-1$, let $\xi$ be uniform on $S^{d-1}$ and define
\[
 U(r)=\E|r e_1+\xi|^{-p}\qquad(r\geq0).
\]
Then $U(r)\leq U(1)$ for every $r\geq0$, with equality only at $r=1$.
\end{lemma}

\begin{proof}
Rotational invariance gives
$U(r)=\E|e_1+r\xi|^{-p}$.  For $0\leq r<1$, the two-coefficient expansion
from the source~\cite[Lemma 13]{CST} is
\begin{equation}\label{eq:hyp}
 U(r)={}_2F_1\left(\frac p2,\frac{p-d+2}{2};\frac d2;r^2\right)
 =\sum_{m=0}^{\infty}
 \frac{(p/2)_m((p-d+2)/2)_m}{(d/2)_m m!}\,r^{2m}.
\end{equation}
Both numerator parameters are positive, so every coefficient is positive.
The series converges at $r=1$ because $p<d-1$; hence $U$ is strictly
increasing on $[0,1]$.

For $r>1$, homogeneity gives
\[
 U(r)=r^{-p}U(1/r).
\]
With $t=r^{-2}$, the right side is $t^{p/2}F(t)$, where $F$ is the
positive-coefficient power series in \eqref{eq:hyp}.  This expression is
strictly increasing in $t\in(0,1)$, so $U$ is strictly decreasing in
$r\in(1,\infty)$.  The unique maximum is therefore at $r=1$.
\end{proof}

\begin{proof}[Proof of Theorem~\ref{thm:main}]
By symmetry take $a=|a_1|=\max_k|a_k|$ and put
$Y=\sum_{k=2}^n a_k\xi_k$.  Condition on $Y$.  Averaging over $\xi_1$ and
rotating the fixed vector $Y$ to the first coordinate gives, by
Lemma~\ref{lem:peak},
\begin{align*}
 \E\left|a\xi_1+Y\right|^{-p}
 &=a^{-p}\E\,U(|Y|/a)\\
 &\leq a^{-p}U(1).
\end{align*}
The law of a pair of independent spherical vectors and rotational invariance
give
\[
 M_{d,2}(p)=2^{p/2}\E|\xi_1+\xi_2|^{-p}
 =2^{p/2}U(1).
\]
Writing $A^2=\sum_k a_k^2$, condition \eqref{eq:large} now yields
\[
 a^{-p}U(1)\leq 2^{p/2}A^{-p}U(1)=M_{d,2}(p)A^{-p},
\]
which proves \eqref{eq:moment}.  Taking both sides to the power $-1/p$
gives the equivalent $L_{-p}$ inequality.

For equality, both steps must be equalities.  Thus $a^2=A^2/2$ and
$|Y|=a$ almost surely.  If the nonzero tail coefficients are $b_j$, then
independence and $\E\langle\xi_i,\xi_j\rangle^2=1/d$ give
\[
 \Var |Y|^2=\frac4d\sum_{i<j}b_i^2b_j^2.
\]
Constancy of $|Y|$ forces at most one nonzero tail coefficient.  Its squared
magnitude is $A^2-a^2=a^2$, proving the equality assertion.  Conversely, two
coefficients of equal magnitude plainly give equality.
\end{proof}

\section{Upgrade audit: the all-small obstruction}

The theorem does not cover $\max a_k^2<A^2/2$.  Two materially different
upgrade routes were pursued.

\paragraph{Three summands in $d=3$.}
When $a,b,c$ are all small, they obey the triangle inequalities.  Put
$\beta=3-p\in(1,2)$ and write
$a=y+z$, $b=x+z$, $c=x+y$.  Explicit spherical convolution reduces the full
three-summand claim, after $x+y+z=1$, to
\begin{equation}\label{eq:scalar}
 (1-x^\beta-y^\beta-z^\beta)(1+x^2+y^2+z^2)^{(3-\beta)/2}
 \leq \beta 2^{(3-\beta)/2}(1-x)(1-y)(1-z).
\end{equation}
The endpoint $\beta=2$ follows from
$xyz\leq(xy+yz+zx)^2/3$; the endpoint $\beta=1$ is trivial.  Finite-grid
checks found no interior violation, but no valid interpolation for the
coupled exponents in \eqref{eq:scalar} was obtained.  It is therefore recorded
only as an unproved route, not as part of the result.

\paragraph{Arbitrary all-small vectors.}
The source's Fourier/H\"older method reduces the missing estimates to sharp
weighted Ball-type inequalities
\[
 s^{p/2}\int_0^\infty |j_{d/2-1}(t)|^s t^{p-1}\,dt
 \leq 2^{p/2}\int_0^\infty |j_{d/2-1}(t)|^2t^{p-1}\,dt,
 \qquad s\geq2.
\]
Here $j_0=J_0$ for $d=2$ and $j_{1/2}(t)=\sin(t)/t$ for $d=3$.  The
integrals converge in exactly the open strips, but equality at $s=2$ makes
nonsharp Gaussian or tail bounds insufficient.  This is the remaining sharp
analytic barrier.

\section{Verification, novelty status, and limitations}

The proof was checked at the exponent endpoints, under scaling, and in the
equality case.  The identity for $\Var|Y|^2$ was independently expanded term
by term.  The script in \texttt{code/check\_upgrade\_reduction.py} checks the
unproved three-summand reduction on finite grids; it is not used as proof.

A bounded search on 13 August 2026 covered the run indexes, arXiv:2206.01026,
the exact open-range sentences, and combinations of ``negative spherical
Khinchin'', ``largest coefficient'', ``half the squared mass'', and
``spherical Riesz potential''.  No explicit statement of this partial theorem
or later full resolution was found.  Novelty confidence is moderate: the key
hypergeometric expansion is in the source and the projection argument is
short, but this application to the two stated open ranges appears unrecorded.

The result is a coefficient-region theorem, not a proof of either complete
conjecture.  Human review should focus on the reciprocity step in
Lemma~\ref{lem:peak}, conditioning at a possibly zero tail, and the equality
case.  The recommended disposition is to retain this as a substantial partial
result and pursue \eqref{eq:scalar} or the weighted Bessel inequality
separately.

\begin{thebibliography}{9}
\bibitem{CST}
G. Chasapis, S. Singh, and T. Tkocz,
\emph{Haagerup's phase transition at polydisc slicing},
arXiv:2206.01026 (2022; version dated 2023).
\url{https://arxiv.org/abs/2206.01026}.
\end{thebibliography}

\end{document}
